\begin{abstract}
Intrinsic self-correction means asking a large language model to critique and revise its own answer without external feedback. It often fails to improve accuracy, because revision can turn correct answers into incorrect ones. One possible fix is confidence gating: revise only when the model's confidence in its initial answer is low. We give a simple analytical model of gated revision. It shows that the gain of any gate is capped by how often the revision step repairs a wrong answer, and that an uninformative signal can never beat the better of never and always revising. We then compare three training-free confidence signals (verbalised confidence, self-consistency and P(True) self-verification) as gates for \modelA{} on GSM8K and HotpotQA, with thresholds tuned on a development split. Self-critique almost never repairs an error (fix rate 0 on GSM8K, 0.04 on HotpotQA), so no gate improves on never revising. Self-consistency, however, identifies wrong answers well (AUROC 0.95 on GSM8K), and its majority answer repairs 47\% of GSM8K errors; used as the revision step, the same gates behave as the analytical model predicts. The results suggest that the bottleneck of self-correction is the revision step, not the decision of when to revise.
\end{abstract}
